\begin{table}[t]
\centering
\small
\caption{接近强消融比较的解释。该表说明每项比较支持什么结论，以及不支持什么结论。}
\label{tab:strong-ablation-interpretation}
\begin{tabularx}{\linewidth}{p{0.27\linewidth}p{0.23\linewidth}X}
\toprule
比较 & 观察结果 & 可支持的解释 \\
\midrule
完整模型 vs. 扁平医学知识 LTR，无图结构 & MRR@10、Recall@10、nDCG@10 均无显著差异 & 扁平检索、语义和生物医学知识特征解释了大部分总体收益；不应声称完整模型显著压倒该消融。 \\
完整模型 vs. 成对图 LTR & $k=10$ 上无显著差异；nDCG@10 是接近显著的正向趋势 & 超图是有用的结构变体，但当前 top-10 证据不足以证明其相对于成对图分解具有清晰总体优势。 \\
完整模型 vs. 移除 MeSH 层级 & 无显著差异；该消融的 Recall@10 和 nDCG@10 略高 & MeSH 层级特征更适合表述为可解释医学约束，而不是已被总体指标单独证明的增益来源。 \\
完整模型 vs. 移除超图扩散/中心性 & 无显著差异；该消融的 MRR@10 和 nDCG@10 略高 & 扩散和中心性特征有助于透明性和案例级解释，但当前测试不能证明它们独立提高总体排序指标。 \\
\bottomrule
\end{tabularx}
\end{table}
